\section{The Manifest}
\label{sec:model}

A manifest is a JSON document with a canonical (RFC~8785) serialization over which a
signature is computed. Table~\ref{tab:top} lists the top-level fields.

\begin{table}[h]
\centering
\small
\caption{Top-level manifest fields (optional fields italicized).}
\label{tab:top}
\begin{tabular}{p{3.7cm}p{3.6cm}p{4.9cm}}
\toprule
Field & Type & Notes \\
\midrule
\texttt{manifest\_id} & UUIDv7~\citep{uuid9562} & time-ordered unique id \\
\texttt{agent\_id} & SPIFFE URI & the agent workload \\
\texttt{version} & string & schema version (\texttt{0.1}) \\
\texttt{issued\_at} & ISO 8601 UTC & issuance time \\
\texttt{expires\_at} & ISO 8601 UTC & $\geq$ 1\,h after, $\leq$ 365\,days \\
\texttt{issuer} & SPIFFE URI & signing authority \\
\texttt{crypto\_profile} & \texttt{standard} $|$ \texttt{post-quantum} & signature suite \\
\texttt{artifacts} & object & the ten artifact bindings (Table~\ref{tab:artifacts}) \\
\texttt{attestation} & object & TEE evidence; required at Level~1+ \\
\emph{\texttt{delegation\_chain}} & list & present when the agent was delegated to \\
\emph{\texttt{hitl\_record}} & object & present when policy mandates human approval \\
\texttt{signature} & object & algorithm, key id, signed-field list, value \\
\emph{\texttt{transparency\_log\_entry}} & object & inclusion proof; required in production \\
\bottomrule
\end{tabular}
\end{table}

\subsection{The ten artifact bindings}

The \texttt{artifacts} object binds each agent-defining surface by hash and metadata,
so any post-approval change to the underlying artifact is detectable.
Table~\ref{tab:artifacts} summarizes the bindings.

\begin{table}[h]
\centering
\small
\caption{The ten artifact bindings and their key bound fields.}
\label{tab:artifacts}
\begin{tabular}{p{0.4cm}p{3.0cm}p{8.6cm}}
\toprule
\# & Binding & Key bound fields \\
\midrule
1 & System prompt & \texttt{hash}, version, classification, safety level \\
2 & Policy bundle & \texttt{hash}, policy language (Cedar/Rego/YAML/composite), enforcement mode, AGT version \\
3 & Tool manifest & \texttt{catalog\_hash}, per-tool entries, dynamic-registration and rug-pull policy \\
4 & Model identity & provider, model id, version, deployment type, \texttt{model\_hash}, attestation type \\
5 & RAG corpus & \texttt{merkle\_root}, document count, ingestion-policy hash, poisoning scan \\
6 & Memory baseline & \texttt{snapshot\_hash}, memory type, drift policy, TTL \\
7 & Decision trace & \texttt{audit\_chain\_root} (a TRACE chain), trace type, signing key id, entry count \\
8 & Delegation chain & per-hop principal, \texttt{scope\_grant}, delegation signature, principal manifest id \\
9 & Supply chain & container/base-image digests, SLSA provenance, SBOM, MCP servers \\
10 & Human-in-the-loop & required flag, approvals, escalation policy \\
\bottomrule
\end{tabular}
\end{table}

Figure~\ref{fig:binding} shows the binding structure: ten heterogeneous artifacts
each reduce to a digest, the digests collect into the \texttt{artifacts} object,
the whole document is serialized with RFC~8785 canonical JSON, and one signature
plus an optional transparency-log inclusion proof close the record. Any change to
any bound artifact changes its digest, changes the canonical bytes, and invalidates
the single signature: this is what makes the record tamper-evident.

\begin{figure}[h]
\centering
\begin{tikzpicture}[
  font=\scriptsize,
  art/.style={draw, rounded corners=2pt, minimum width=2.55cm, minimum height=0.42cm,
    fill=blue!5, align=left, inner sep=2pt},
  core/.style={draw, thick, rounded corners=3pt, minimum width=2.7cm, minimum height=0.9cm,
    fill=blue!12, align=center},
  sig/.style={draw, thick, rounded corners=3pt, minimum width=2.7cm, minimum height=0.7cm,
    fill=green!12, align=center},
  tl/.style={draw, rounded corners=3pt, minimum width=2.7cm, minimum height=0.65cm,
    fill=orange!12, align=center},
  ln/.style={-{Stealth[length=1.6mm]}, gray!70},
]
% Left column: ten artifact bindings
\node[art] (a1)  {1 System prompt};
\node[art] (a2)  [below=2pt of a1] {2 Policy bundle};
\node[art] (a3)  [below=2pt of a2] {3 Tool manifest};
\node[art] (a4)  [below=2pt of a3] {4 Model identity};
\node[art] (a5)  [below=2pt of a4] {5 RAG corpus};
\node[art] (a6)  [below=2pt of a5] {6 Memory baseline};
\node[art] (a7)  [below=2pt of a6] {7 Decision trace};
\node[art] (a8)  [below=2pt of a7] {8 Delegation chain};
\node[art] (a9)  [below=2pt of a8] {9 Supply chain};
\node[art] (a10) [below=2pt of a9] {10 Human-in-the-loop};

% Middle: the artifacts object (collected digests)
\node[core] (obj) [right=2.1cm of a5, yshift=-9pt] {\texttt{artifacts}\\object\\(10 digests)};

% edges from each artifact to the object, labeled as a digest
\foreach \n in {a1,a2,a3,a4,a5,a6,a7,a8,a9,a10} {
  \draw[ln] (\n.east) -- (obj.west);
}
\node[gray!80] at ($(a5.east)!0.5!(obj.west) + (0,0.95)$) {\itshape hash / Merkle root};

% Right: canonicalize -> sign -> transparency log
\node[core] (canon) [right=1.3cm of obj] {RFC~8785\\canonical\\JSON bytes};
\node[sig]  (sign)  [below=0.5cm of canon] {Ed25519 / ML-DSA-65\\signature};
\node[tl]   (tlog)  [below=0.45cm of sign] {transparency log\\inclusion proof};

\draw[ln] (obj.east) -- (canon.west);
\draw[ln] (canon.south) -- (sign.north);
\draw[ln] (sign.south) -- (tlog.north);

% background box around the signed record
\begin{scope}[on background layer]
  \node[draw, dashed, rounded corners, fit=(canon)(sign)(tlog), inner sep=6pt,
        label={[font=\scriptsize\itshape]above:signed, tamper-evident record}] {};
\end{scope}
\end{tikzpicture}
\caption{Ten agent-defining artifacts each reduce to a digest, collect into the
\texttt{artifacts} object, canonicalize under RFC~8785, and are closed by one
signature and an optional transparency-log inclusion proof. Any post-approval change
to any artifact changes its digest and invalidates the single signature.}
\label{fig:binding}
\end{figure}

The binding pattern is uniform: each artifact contributes a digest (a SHA-256 hash, a
Merkle root for collections, or an audit-chain root for the decision trace) into the
signed manifest, so verifying the signature plus recomputing the artifact's digest
detects any divergence. The model-identity binding distinguishes a \emph{hash-bound}
model (weights digest verifiable) from a \emph{provider-asserted} one (a hosted API,
where the provider's assertion is the strongest available evidence). The
decision-trace binding is the join point with the runtime stack: its
\texttt{audit\_chain\_root} is the root of the TRACE audit chain cMCP
produces~\citep{trace2026,cmcp2026}.
